\section{Introduction}

Signal reconstruction from corrupted or incomplete measurements remains a critical challenge in statistical data analysis, particularly within audio processing where artifacts such as periodic seam-wrap crackle degrade perceptual quality \cite{stoller2018}. While deep neural networks have revolutionized feature extraction for tasks like speech recognition and source separation, their training efficacy is often hindered by sensitivity to hyperparameter choices including model architecture and loss functions \cite{jaderberg2017}. This fragility necessitates robust meta-learning frameworks capable of navigating vast search spaces without relying on clean ground-truth data. We address this gap by proposing a hybrid reinforcement learning (RL) and genetic algorithm (GA) approach that simultaneously optimizes policy parameters, network topology, and architectural components via population-based training \cite{jaderberg2017}. Our method integrates differentiable digital signal processing primitives to ensure physical consistency in the generated waveforms \cite{engel2020}, thereby bridging the divide between abstract optimization objectives and acoustic reality.

Existing literature has established foundational techniques for learning from corrupted data, such as Noise2Noise which demonstrates that restoration is possible using only noisy observations through statistical reasoning \cite{lehtinen2018}. Similarly, evolutionary strategies have been successfully applied to discover neural network topologies for image classification tasks, challenging the necessity of manual expert design \cite{real2019}. However, these approaches often treat architecture search and policy optimization as decoupled processes. Recent advances in Regularized Evolution suggest that direct evolution can outperform gradient-based methods when combined with efficient parameter sharing mechanisms to reduce computational overhead \cite{pham2018;elsken2019}. Furthermore, the integration of sparse gating via Mixture-of-Experts layers offers a pathway to increase model capacity without proportional increases in parameters, addressing scalability concerns inherent in large-scale audio synthesis tasks \cite{shazeer2017}. Despite these advancements, current Deep Reinforcement Learning algorithms still struggle with temporal credit assignment under sparse rewards and lack effective exploration strategies required for complex control domains like audio denoising \cite{khadka2018;schulman2017}.

To overcome these limitations, we introduce a unified framework that couples Proximal Policy Optimization (PPO) with genetic algorithms to jointly refine the policy network and its underlying architecture. This hybrid approach leverages population-based training principles to maintain diversity within the search space, preventing premature convergence often observed in standard evolutionary methods \cite{jaderberg2017}. By embedding differentiable DSP operations directly into the loss function, our model learns to respect the periodic boundary conditions essential for wavetable synthesis, effectively mitigating seam artifacts. We distinguish between used literature that informs our specific methodological choices and screened works providing broader context on audio source separation in time domains \cite{stoller2018;purwins2019}. While comprehensive benchmarks are currently underway to validate these claims across diverse acoustic environments, preliminary analysis suggests significant potential for improving residual metrics compared to historical baselines. This work contributes a novel synthesis of evolutionary computation and differentiable signal processing aimed at solving the persistent problem of periodic noise removal in synthesized audio signals.
